% chapters/13/contoh.tex
\section{Dua Contoh Analisis}
\label{sec:contoh-bab-13}

\subsection{Contoh Pertama: Membaca Tangkapan Layar Level}

\begin{contohbox}
Perhatikan kembali Gambar~\ref{fig:editor-fisika}. Tangkapan layar itu
memperlihatkan sebuah level yang sedang disusun, dengan bentuk batas benda
ditampilkan oleh panel \istilah{Physics Debug}.
\end{contohbox}

Tabel~\ref{tab:baca-level} memuat pembacaan bagian-bagian tangkapan layar itu.

\begin{table}[htbp]
\centering
\caption{Pembacaan bagian-bagian tangkapan layar level}
\label{tab:baca-level}
{\setlength{\tabcolsep}{4pt}\footnotesize%
\begin{tabularx}{\textwidth}{@{}L{4.2cm} Y@{}}
\toprule
\headrow \textbf{Yang tampak} & \textbf{Keterangan} \\
\midrule
Balok abu-abu tanpa tekstur, tangga, dan bongkahan &
Level yang disusun secara \istilah{greyboxing}. Belum ada aset yang dipasang,
sebab yang diuji masih mekaniknya. \\
\addlinespace
Kotak berwarna kuning dengan garis tepi jingga &
Bentuk batas sebuah benda, yaitu \istilah{Collider}, yang ditampilkan oleh
\istilah{Physics Debug}. Bentuknya sederhana, yaitu berupa
kotak, walaupun gambarnya lebih rumit. \\
\addlinespace
Tanda panah berwarna pada benda itu &
\istilah{Gizmo}, yaitu tanda untuk memindahkan dan memutar benda pada panel
\istilah{Scene}. Warnanya menunjukkan sumbu: merah untuk mendatar, hijau untuk
tegak. \\
\addlinespace
Panel kecil bertulisan \istilah{Physics Debug} pada sudut kanan bawah &
Panel pengaturan tampilan fisika, dengan pilihan \istilah{Collision Geometry}
dan \istilah{Mouse Select}. \\
\bottomrule
\end{tabularx}}
\end{table}

Penyusunan level dan penampilan bentuk batasnya dapat
dikerjakan bersamaan. Kebiasaan menampilkan bentuk batas sejak awal menyusun
level akan menghemat banyak waktu, sebab masalah fisika dapat ditemukan
sebelum levelnya menjadi panjang.

\subsection{Contoh Kedua: Menelusuri Tiga Gejala sampai Penyebabnya}

\begin{contohbox}
Seorang mahasiswa melaporkan tiga gejala pada levelnya: pemain bergetar ketika
berjalan, pemain kadang menembus lantai ketika jatuh dari tempat tinggi, dan
sebuah bongkahan tidak dapat dipijak sama sekali. Ketiganya ditelusuri menurut
urutan pada Bagian~\ref{sec:masalah-fisika}.
\end{contohbox}

Hasil penelusurannya diperlihatkan pada Tabel~\ref{tab:telusur-fisika}.

\begin{table}[htbp]
\centering
\caption{Hasil penelusuran tiga gejala fisika}
\label{tab:telusur-fisika}
{\setlength{\tabcolsep}{4pt}\footnotesize%
\begin{tabularx}{\textwidth}{@{}L{2.6cm} L{3.4cm} L{3.2cm} Y@{}}
\toprule
\headrow \textbf{Gejala} & \textbf{Hasil pemeriksaan} & \textbf{Penyebab} & \textbf{Perbaikan} \\
\midrule
Pemain bergetar ketika berjalan &
Skrip menggerakkan \code{transform.position} pada \code{Update}. &
Perubahan letak terjadi di antara dua langkah fisika, sehingga kamera
menampilkan letak yang berbeda dari letak yang dihitung fisika. &
Gerakkan melalui \istilah{Rigidbody} pada \code{FixedUpdate}, dan aktifkan
\istilah{Interpolate}. \\
\addlinespace
Pemain menembus lantai ketika jatuh dari tempat tinggi &
Kecepatan jatuh mencapai 40 satuan per detik, sedangkan tebal lantai 0,2
satuan. &
Pada satu langkah fisika, perpindahannya sekitar 0,67 satuan, yaitu lebih tebal
daripada lantainya. &
Ubah \istilah{Collision Detection} menjadi \istilah{Continuous}, atau
tebalkan lantainya. \\
\addlinespace
Bongkahan tidak dapat dipijak &
Bongkahan tidak memiliki \istilah{Collider}, sebab hanya diberi komponen
gambar. &
Tanpa \istilah{Collider}, benda tidak memiliki bentuk batas, sehingga mesin
tidak memeriksanya sama sekali. &
Tambahkan \istilah{Box Collider 2D} pada bongkahan itu. \\
\bottomrule
\end{tabularx}}
\end{table}

Dua dari ketiga gejala itu bukan kesalahan kode, melainkan kesalahan
pengaturan. Perbandingan itu penting untuk diperhatikan: pada blok Python,
seluruh masalah berada pada kode Anda, sedangkan pada Unity, pengaturan
komponen justru menjadi tempat pertama yang perlu diperiksa.
